%%%%%%%%%%%%%%%%%%%%%%%%%%%%%%%%%%%%%%%%%%%%%%%%%%%%%%%%%%%%%%%%%%%%%%
\section{Data}\label{sec:data}
%%%%%%%%%%%%%%%%%%%%%%%%%%%%%%%%%%%%%%%%%%%%%%%%%%%%%%%%%%%%%%%%%%%%%%

The model is built entirely from publicly and freely available data. Four
groups of inputs are combined: (i)~Brazilian subnational statistics on soy
production, processing, domestic use and trade; (ii)~national commodity
balances and feed parameters; (iii)~administrative and environmental
geographies; and (iv)~the global input--output backends and the independent
validation benchmark. Unless otherwise stated, the series are annual, cover the period 2001–2022 and are broken down by Brazil's approximately 5,570 municipalities
(\emph{munic\'ipios}). Where a source exists for only a single reference
year, that year is stated and the layer is held constant.
Table~\ref{tab:data} summarises every input, and all access links are collected
in Appendix~\ref{app:links}. A downloader for the machine-retrievable
sources is included with the code.

% ==================================================================
\subsection{Brazilian subnational statistics}\label{subsec:brazil}
% ==================================================================

\textbf{Production and population (IBGE):} Municipal soybean production and the planted and harvested area data come from the SIDRA system (Table 1612), obtained via its public API. The resident population (Table 6579; 2022 Census) serves as a proxy for household soy oil consumption.

\textbf{Livestock (IBGE):} Municipal herds for ten animal categories (Table 3939) and milked cow counts (Table 94) provide the feed demand base. Feedlot cattle, last reported in the 2006 agricultural census (the 2017 census dropped the confinement question), are carried forward from the 2006 municipal pattern using the national ABIEC/Athenagro confinement series (see Methods).

\textbf{Processing and biodiesel capacity (ABIOVE, ANP):} Soy crushing and
refining capacity comes from two manually downloaded ABIOVE products: the
consolidated state-level installed capacities and the annual plant roster
surveys (\emph{Pesquisa de Capacidade}), whose changing sheet layouts are
handled by a per-year dispatcher in the code. Refining and bottling
capacities come from the refining sheets of the same products. Years for which surveys are missing from the ABIOVE archive (2000–2002, 2016–2017 and 2021) are replaced with the previous available survey. Biodiesel plant capacity and the regional soy feedstock share are taken from the ANP Statistical Yearbook.

\textbf{Household consumption and storage (IBGE):} State-level per capita soy oil acquisition from the Household Budget Survey (POF 2017–18, the most recent) is used to distribute national food oil use and is held constant across years. Grain storage capacity, which allocates stock changes, is aggregated from the georeferenced warehouse registry underlying the IBGE Stocks Survey (2014 snapshot, held constant).

\textbf{Municipal trade (COMEX/MDIC):} Municipality-level exports and imports are sourced from the COMEX Stat customs database and filtered by the HS4 soy codes 1201 (beans), 1507 (oil) and 2304 (cake).

% ==================================================================
\subsection{National balances and feed parameters (FAO)}\label{subsec:fao}
% ==================================================================

\textbf{Commodity balances:} The national supply-use balances against which all municipal totals are reconciled are the FAO Food Balance Sheets for Brazil, items 2555 (soybeans), 2571 (soybean oil) and 2590 (soybean cake), taken from the FAOSTAT bulk download.

\textbf{Bilateral trade:} Re-export disentangling draws on the FAOSTAT
Detailed Trade Matrix (items 236, 237 and 238 for beans, oil and cake).

\textbf{Livestock systems and feed rates:} Herds are categorised by production system using the Gridded Livestock of the World rasters (GLW3, reference year 2010) \cite{gilbert2018glw} and the GLEAM ruminant system map. They are then converted to soy feed demand using GLEAM per-animal intake rates \cite{fao2018gleam}.. These are static layers: the resulting system shares are held constant while herd sizes remain annual (see Methods).

% ==================================================================
\subsection{Geographies}\label{subsec:geo}
% ==================================================================

\textbf{Boundaries:} Municipal polygons are IBGE administrative meshes
retrieved with the \texttt{geobr} package \cite{pereira2019geobr}; state
outlines for mapping use GADM v3.6.

\textbf{Biomes and frontier:} The footprint origin is attributed to biomes using the IBGE 1:250{,}000 biome layer. The agricultural frontier is the Amazon biome plus MATOPIBA, operationalised as the Cerrado biome within Maranh\~ao, Tocantins, Piau\'i and Bahia, closely matching the official Embrapa delineation \cite{miranda2014matopiba}.

\textbf{Grid refinement (optional):} A 30\,m visual refinement uses
MapBiomas annual soy tiles \cite{souza2020mapbiomas}, exported via Google
Earth Engine. This refinement is illustrative only and does not affect the load-bearing capacity of the results.

\textbf{Transport network:} The reported model routes flows based on straight-line distances between municipal centres (see Methods). Four multimodal layers — OpenStreetMap roads, DNIT waterways, ANTT rail lines and cargo, and ANTAQ ports and port cargo — support an experimental cost-based transport variant stored in the repository and used for descriptive analyses. These layers do not affect the reported results.

% ==================================================================
\subsection{Global models and benchmark}\label{subsec:mrio-data}
% ==================================================================

\textbf{FABIO.} Subnational flows are incorporated into the Food and Agriculture
Biomass Input--Output model \cite{bruckner2019fabio}. From 2010 onwards the
version~2 build is used (2010--2023; 22{,}263 processes, 181 regions
$\times$ 123 items, with 123 environmental extensions). As the v2 build did not exist before 2010, the years 2001–2009 use the corresponding inputs of the v1.1 build: bilateral trade details, land use extensions, feedstock optimisation results and v1 balanced trade and commodity balance sheets. Version bridging and the quantified v1.1/v2 seam are described in the 'Methods' section.

\textbf{EXIOBASE.} The non-soy background economy is EXIOBASE~3
\cite{stadler2018exiobase} in the product-by-product transformation
(2000--2022; 49 regions $\times$ 200 products), hybridised with FABIO for
the non-food side of the footprint.

\textbf{Trase (validation only).} Modelled origin-to-importer flows are
validated against the Trase Brazil soy composite v2.6.1 (2004--2022), which
links producing municipalities to importing countries from per-shipment
customs and logistics records \cite{godar2015sei,escobar2020trase}. No
Trase data enter the model itself.

% ==================================================================
\subsection{Coverage and summary}\label{subsec:coverage}
% ==================================================================

The Footprint series is reported for the period 2001–2022. FABIO v2 covers the period from 2010 onwards, while the FABIO v1.1 build provides the backend for the period from 2001 to 2009. The year 2000 is excluded pending the resolution of a demand-side allocation issue in that year's build (see Methods). The Trase benchmark supports validation for 2004–2022.

\begin{table}[htbp]
\small
\caption{Input datasets, by role. All are publicly and freely available;
access links are collected in Appendix~\ref{app:links}. ``Munis'' $=$
Brazilian municipalities; ``exp.'' $=$ experimental variant only.}\label{tab:data}
\begin{tabular}{@{}p{3.5cm}p{2.3cm}p{3.0cm}p{2.6cm}@{}}
\toprule
Dataset & Provider & Role & Coverage \\
\midrule
Soy production, area (t.1612)      & IBGE (SIDRA)        & production            & munis, annual \\
Population (t.6579/9514)           & IBGE (SIDRA)        & food-use proxy        & munis, annual \\
Herds (t.3939), milk cows (t.94)   & IBGE (SIDRA)        & feed demand           & munis, annual \\
Feedlot cattle (t.919)             & IBGE (Census)       & feed demand           & munis, 2006 pattern \\
Crush/refining capacity            & ABIOVE              & processing            & munis, annual \\
Biodiesel capacity                 & ANP                 & other (biodiesel) use & munis, annual \\
Soy-oil acquisition (POF)          & IBGE                & food-oil allocation   & states, 2017--18 \\
Grain storage (warehouse registry) & IBGE                & stock allocation      & munis, 2014 static \\
Commodity balance (FBS)            & FAO (FAOSTAT)       & national reconciliation & Brazil, annual \\
Gridded livestock (GLW3)           & FAO (Dataverse)     & system split          & 10\,km, 2010 \\
Production systems, feed (GLEAM)    & FAO                 & feed rates            & 10\,km, static \\
Municipal trade                    & COMEX/MDIC          & exports/imports       & munis, annual \\
Road network                       & OSM (Geofabrik)     & transport (exp.)      & national \\
Waterways                          & DNIT                & transport (exp.)      & national \\
Rail lines, stations, cargo        & ANTT                & transport (exp.)      & national, 2006--21 \\
Ports, port cargo                  & ANTAQ               & transport (exp.)      & national, 2017--22 \\
Municipal boundaries               & IBGE (geobr)        & geometry              & munis \\
State boundaries                   & GADM v3.6           & mapping               & states \\
Biomes 1:250k                      & IBGE                & origin attribution    & national \\
Soy land tiles                     & MapBiomas (GEE)     & grid refinement       & 30\,m, annual \\
FABIO v2 (MRIO)                    & fineprint           & supply-chain nesting, 2010-- & global, 2010--23 \\
FABIO v1.1 (MRIO)                  & fineprint           & supply-chain nesting, pre-2010 & global, 1986--2013 \\
EXIOBASE 3 pxp                     & EXIOBASE            & non-soy background    & global, 2000--22 \\
Detailed trade matrix              & FAO (FAOSTAT)       & re-exports/trade      & Brazil, annual \\
Trase composite v2.6.1             & Trase               & validation benchmark  & munis, 2004--22 \\
\botrule
\end{tabular}
\end{table}

% ==================================================================
\subsection{Data availability}\label{subsec:data-availability}
% ==================================================================
All input data are publicly and freely available from the providers listed
in Appendix~\ref{app:links}. They are not redistributed with the code
because of differing licences; a download script covers the
machine-retrievable sources, and the datasets requiring manual retrieval
(ABIOVE, ANP, POF, MapBiomas tiles) are documented step by step in the
repository. The intermediate and final model outputs underlying the figures
will be deposited on Zenodo upon publication.

% ==================================================================
\subsection{Code availability}\label{subsec:code-availability}
% ==================================================================
The complete modelling pipeline (data preparation, transport allocation,
FABIO/EXIOBASE nesting, footprint accounting, validation and figures) is
written in R and Python and released under the GNU GPL~v3 at
\texttt{<repository URL --- TODO>}. The results reported here have no
proprietary software dependency. The experimental multimodal-transport
module included in the repository (not used for the reported results)
additionally requires a GAMS licence.
